\begin{topic}[id=fiveg_m2m,
             difficulty={Mittel},
             type={Systemanalyse + Einordnung},
             prerequisites={Grundlagen Mobilfunk und Rechnernetze; Interesse an industrieller Kommunikation},
             practice={optional},
             owner={Dozententeam},
             updated={2026-04-09},
             tags={ 5G, M2M, URLLC, eMBB, mMTC, Slicing, Edge, Private 5G }]{5G \& M2M}

\TopicDescription{5G verspricht sehr niedrige Latenzen (URLLC), hohe Bandbreiten (eMBB) und Massengeräte (mMTC). Für industrielle Anwendungen sind Netzslicing, Edge-Integration und Private-5G-Campusnetze relevant. Wir klären die real erreichbaren Eigenschaften, Abhängigkeiten zu Endgeräten/Netzinfrastruktur und den Betriebsaufwand. Ziel: eine nüchterne Einordnung, wann 5G-M2M Vorteile bringt und wann Alternativen (Wi-Fi 6/7, LTE-M, NB-IoT) genügen.}

\TopicLeitfragen{\item Welche 5G-Features sind in der Praxis heute nutzbar?
\item Wann lohnt Private 5G gegenüber Wi-Fi + VLAN/QoS?
\item Wie integriert man Edge-Computing sinnvoll?
}
\TopicScope{\item Keine Frequenzvergaben/Regulatorik im Detail.
\item Keine Hardware-Benchmarks.
}
\TopicPractice{Optionale Praxisidee: Architektur-Skizze mit Slicing/Edge-Integration für eine Beispielanlage.}

\TopicKeywords{5G, M2M, URLLC, eMBB, mMTC, Slicing, Edge, Private 5G}

\nocite{3gppOverview,etsi5g}
\end{topic}
